% Presentación Beamer autocontenida (16:9). Compilar con pdfLaTeX dos veces.
\documentclass[aspectratio=169,10pt]{beamer}

\usepackage[utf8]{inputenc}
\usepackage[T1]{fontenc}
\usepackage[spanish,es-tabla,es-nodecimaldot]{babel}
\usepackage{lmodern}
\usepackage{amsmath,amssymb}
\usepackage{array,tabularx}
\usepackage{tikz}
\usetikzlibrary{arrows.meta,positioning}

\usetheme{Madrid}
\usecolortheme{default}

\definecolor{utecazul}{RGB}{11,79,108}
\definecolor{uteccian}{RGB}{0,145,189}
\definecolor{utecgris}{RGB}{242,246,250}
\definecolor{utecoscuro}{RGB}{32,42,54}

\setbeamercolor{structure}{fg=utecazul}
\setbeamercolor{palette primary}{bg=utecazul,fg=white}
\setbeamercolor{palette secondary}{bg=uteccian,fg=white}
\setbeamercolor{palette tertiary}{bg=utecoscuro,fg=white}
\setbeamercolor{title}{bg=utecazul,fg=white}
\setbeamercolor{frametitle}{bg=utecazul,fg=white}
\setbeamercolor{block title}{bg=utecazul,fg=white}
\setbeamercolor{block body}{bg=utecgris,fg=black}
\setbeamercolor{alertblock title}{bg=uteccian,fg=white}
\setbeamercolor{alertblock body}{bg=utecgris,fg=black}

\setbeamertemplate{navigation symbols}{}
\setbeamertemplate{footline}{%
  \leavevmode%
  \hbox{%
    \begin{beamercolorbox}[wd=.38\paperwidth,ht=2.3ex,dp=1ex,leftskip=2ex]{author in head/foot}%
      \usebeamerfont{author in head/foot}UTEC \textbullet\ Visión por Computadora
    \end{beamercolorbox}%
    \begin{beamercolorbox}[wd=.44\paperwidth,ht=2.3ex,dp=1ex,center]{title in head/foot}%
      \usebeamerfont{title in head/foot}TurtleBot 4 Lite (ROS~2 Jazzy)
    \end{beamercolorbox}%
    \begin{beamercolorbox}[wd=.18\paperwidth,ht=2.3ex,dp=1ex,rightskip=2ex,right]{date in head/foot}%
      \usebeamerfont{date in head/foot}\insertframenumber{} / \inserttotalframenumber
    \end{beamercolorbox}%
  }%
  \vskip0pt%
}

\newcolumntype{Y}{>{\raggedright\arraybackslash}X}
\newcolumntype{P}[1]{>{\raggedright\arraybackslash}p{#1}}
\renewcommand{\arraystretch}{1.15}

\title[TurtleBot 4 Lite: Visión y Navegación]{Navegación autónoma, control gestual y reconocimiento facial en TurtleBot 4 Lite}
\subtitle{Sustentación técnica del proyecto final (8--10 minutos)}
\author[Equipo UTEC]{[Integrante 1] \and [Integrante 2] \and [Integrante 3] \and [Integrante 4]\\[0.3em]\footnotesize Docente: Rensso Victor Hugo Mora Colque}
\institute[UTEC]{Universidad de Ingeniería y Tecnología (UTEC)}
\date{17 de noviembre de 2026}

\begin{document}

% -------------------------------------------------------------------------
\begin{frame}[plain]
  \titlepage
\end{frame}

% -------------------------------------------------------------------------
\begin{frame}{Reto de competición y enfoque técnico}
  \small
  \begin{columns}[T,onlytextwidth]
    \column{0.47\textwidth}
    \begin{block}{Task A: Navegación autónoma}
      \begin{itemize}
        \item Ingresar el \textbf{nombre objetivo} antes de salir de A.
        \item Recorrer circuito desconocido evitando obstáculos y siguiendo \textbf{flechas azules}.
        \item Confirmar meta B, ubicar la foto del objetivo y leer únicamente su \textbf{código QR}.
        \item Retornar guiado por \textbf{flechas rojas} a A y mostrar la palabra QR.
      \end{itemize}
    \end{block}

    \column{0.47\textwidth}
    \begin{block}{Task B: Teleoperación por gestos}
      \begin{itemize}
        \item Conducción exclusiva mediante \textbf{gestos visuales humanos} en video (sin teclado ni giros autónomos).
        \item Ir de A a B, identificar la fotografía de B con la cámara del robot, regresar a A y anunciar el nombre.
      \end{itemize}
    \end{block}
    \vspace{0.45em}
    \begin{alertblock}{Principio rector}
      Circuito y fotografía se revelan en pista: \textbf{cero rutas memorizadas} y \textbf{parada segura} ante datos ausentes o caducados.
    \end{alertblock}
  \end{columns}
\end{frame}

% -------------------------------------------------------------------------
\begin{frame}{1. Arquitectura general del sistema}
  \centering
  \begin{tikzpicture}[
    node distance=0.55cm and 1.25cm,
    caja/.style={draw=utecazul, thick, rounded corners=3pt, fill=utecgris, text width=2.95cm, inner sep=4pt, align=center, minimum height=0.95cm, font=\scriptsize},
    etq/.style={font=\tiny\color{utecoscuro}, align=center}
  ]
    \node[caja] (cam) {\textbf{Host CPU: Visión}\\Cámara OAK-D (5 FPS)\\YuNet + SFace\\Detector Canny/HSV};
    \node[caja, right=1.25cm of cam] (sup) {\textbf{Host: Supervisor}\\Objetivo e histéresis\\Máquina de estados\\Panel de verificación};
    \node[caja, right=1.25cm of sup] (ros) {\textbf{Docker ROS~2 Jazzy}\\Puerta única de salida\\Evasión LiDAR + TF\\Giros por odometría};

    \node[caja, below=0.55cm of cam] (web) {\textbf{Host CPU: Webcam}\\MediaPipe Hand (21 pts)\\Reglas en píxeles\\Confirmación $0.4\ \mathrm{s}$};
    \node[caja, right=1.25cm of web] (pte) {\textbf{Docker: Puente Gestos}\\Vigilante $20\ \mathrm{Hz}$\\Caducidad $0.5\ \mathrm{s}$\\Reloj monotónico};
    \node[caja, right=1.25cm of pte, fill=utecazul!12] (bot) {\textbf{TurtleBot 4 Lite}\\Create~3 (\texttt{/cmd\_vel})\\RPLIDAR (\texttt{/scan})\\Odometría y OAK-D};

    \draw[-{Latex}, thick, utecazul] (cam) -- node[above, etq]{eventos} (sup);
    \draw[-{Latex}, thick, utecazul] (sup) -- node[above, etq]{JSON + latido} (ros);
    \draw[-{Latex}, thick, utecazul] (ros) -- node[right, etq]{velocidad} (bot);
    \draw[-{Latex}, thick, uteccian] (web) -- node[above, etq]{JSON stdin} (pte);
    \draw[-{Latex}, thick, uteccian] (pte) -- node[above, etq]{TwistStamped} (bot);
  \end{tikzpicture}

  \vspace{0.35em}
  \footnotesize
  \begin{columns}[T,onlytextwidth]
    \column{0.47\textwidth}
    \begin{block}{Separación percepción vs. control}
      Permite probar algoritmos en CPU sin acceso físico al robot y aisla la
      versión de Python del host respecto a ROS~2 Jazzy (Ubuntu~24.04).
    \end{block}
    \column{0.47\textwidth}
    \begin{block}{Canal de video aislado}
      La cámara del robot viaja por tubería separada (máx. $5\ \mathrm{FPS}$)
      para que fotogramas pesados nunca bloqueen órdenes de parada.
    \end{block}
  \end{columns}
\end{frame}

% -------------------------------------------------------------------------
\begin{frame}{2. Implementación del reconocimiento facial}
  \footnotesize
  \begin{columns}[T,onlytextwidth]
    \column{0.495\textwidth}
    \begin{block}{Flujo de inferencia en CPU}
      \begin{enumerate}
        \item \textbf{Reescalado:} lado máximo $960\ \mathrm{px}$.
        \item \textbf{YuNet (2023mar):} detecta rostro y 5 puntos clave; exige \textbf{exactamente 1 cara útil} (rechaza 0 o múltiples).
        \item \textbf{Alineación + SFace (2021dec):} extrae descriptor $f(x)\in\mathbb{R}^{128}$ normalizado:
        \[
          z(x)=\frac{f(x)}{\lVert f(x)\rVert_2}, \qquad c(x,r)=z(x)^\top z(r).
        \]
        \item \textbf{Puntuación por identidad:} media de las $k_i=\min(2,n_i)$ mejores referencias.
      \end{enumerate}
    \end{block}

    \column{0.445\textwidth}
    \begin{block}{Criterio dual con rechazo explícito}
      \[
        \widehat y =
        \begin{cases}
          i^\star, & S_{i^\star}(x)\ge\tau\ \text{y}\ \Delta\ge\mu,\\
          \text{NO\_CONCL.}, & \text{en otro caso}.
        \end{cases}
      \]
      \begin{itemize}
        \item Calibración separada de galería y evaluación reservada (manifiesto SHA-256).
        \item En vivo: selección previa del objetivo y confirmación tras \textbf{3 fotogramas nuevos}.
      \end{itemize}
    \end{block}
  \end{columns}
\end{frame}

% -------------------------------------------------------------------------
\begin{frame}{3. Reconocimiento de gestos y teleoperación (Task B)}
  \small
  \begin{columns}[T,onlytextwidth]
    \column{0.47\textwidth}
    \begin{block}{Clasificador geométrico sobre MediaPipe}
      \begin{itemize}
        \item \textbf{21 puntos clave} de la mano proyectados a coordenadas en píxeles (preserva relación de aspecto).
        \item Comparación de extensión/flexión de los 4 dedos largos (bandas ambiguas detienen).
        \item Imagen sin espejo visual; inversión única del eje horizontal para alinear derecha/izquierda del operador con el robot.
      \end{itemize}
    \end{block}

    \column{0.47\textwidth}
    \begin{block}{Tabla de órdenes y seguridades}
      \footnotesize
      \begin{tabularx}{\textwidth}{@{}Ycc@{}}
        \hline
        \textbf{Gesto / estado} & \textbf{Lin. (m/s)} & \textbf{Ang. (rad/s)}\\
        \hline
        Índice arriba & $0.306$ & $0.00$\\
        Índice izquierda & $0.00$ & $+1.90$\\
        Índice derecha & $0.00$ & $-1.90$\\
        Palma / ausencia / $>1$ mano & $0.00$ & $0.00$\\
        \hline
      \end{tabularx}
      \vspace{0.3em}
      \begin{itemize}
        \item Histéresis de \textbf{0.4 s} para avanzar/girar; cambio de sentido pasa por parada.
        \item Vigilante a $20\ \mathrm{Hz}$: caducidad $>0.5\ \mathrm{s}$ detiene.
      \end{itemize}
    \end{block}
  \end{columns}
\end{frame}

% -------------------------------------------------------------------------
\begin{frame}{4. Estrategia de navegación autónoma (Task A)}
  \small
  \begin{columns}[T,onlytextwidth]
    \column{0.47\textwidth}
    \begin{block}{Evasión LiDAR conservadora}
      \begin{itemize}
        \item Sectores: frontal ($\pm 35^\circ$) e izquierdo/derecho.
        \item \textbf{Cobertura mínima $30\%$} ($\ge 3$ rayos finitos) y $\le 10\%$ errores; sector sin retorno suficiente = \textbf{parada}.
        \item Velocidad: $0.10\ \mathrm{m/s}$, giro $0.40\ \mathrm{rad/s}$; precaución $0.45\ \mathrm{m}$, parada frontal $0.32\ \mathrm{m}$, lateral $0.28\ \mathrm{m}$.
      \end{itemize}
    \end{block}

    \column{0.47\textwidth}
    \begin{block}{Señales de giro y maniobra de $90^\circ$}
      \begin{itemize}
        \item Canny + elipses ($\ge 55\ \mathrm{px}$) + filtro HSV por tramo (\textbf{azul} ida / \textbf{rojo} vuelta) + geometría de flecha.
        \item Confirmación en \textbf{3 frames nuevos}; la flecha prevalece sobre el lado de evasión ante un muro.
        \item Giro $90^\circ \pm 5^\circ$ por odometría (cruce $\pm\pi$, timeout $12\ \mathrm{s}$, rearme $1.5\ \mathrm{s}$ y $0.4\ \mathrm{m}$).
      \end{itemize}
    \end{block}
  \end{columns}
  \vspace{0.35em}
  \begin{alertblock}{Pendiente de integración final}
    \footnotesize Acoplar detección de meta A/B por etapa, barrido en B, lectura exclusiva del QR del rostro objetivo y retorno completo.
  \end{alertblock}
\end{frame}

% -------------------------------------------------------------------------
\begin{frame}{5. Integración de sensores y puerta única de seguridad}
  \begin{table}
    \centering
    \footnotesize
    \begin{tabularx}{\textwidth}{@{}P{2.6cm}P{4.1cm}Y@{}}
      \hline
      \textbf{Sensor / canal} & \textbf{Validación previa al movimiento} & \textbf{Condición de parada segura}\\
      \hline
      \textbf{RPLIDAR (\texttt{/scan})} &
      Transformación TF \texttt{base\_link} $\rightarrow$ \texttt{rplidar\_link} y cobertura por sector &
      Antigüedad $>0.6\ \mathrm{s}$, falta de TF o cobertura $<30\%$\\
      \textbf{Odometría (\texttt{/odom})} &
      Finitud, norma unitaria del cuaternión y consistencia de \texttt{frame\_id} &
      Antigüedad $>0.6\ \mathrm{s}$, cambio de marco o pose inválida\\
      \textbf{Cámara OAK-D} &
      QoS de sensores, validación de dimensiones y \texttt{step} de fila &
      Antigüedad $>1.0\ \mathrm{s}$ o fotograma repetido en confirmación\\
      \textbf{Latido del host} &
      Secuencia estrictamente creciente y marca monotónica compartida &
      Expiración $>0.5\ \mathrm{s}$, mensaje malformado o cierre de enlace\\
      \hline
    \end{tabularx}
  \end{table}

  \small
  \begin{block}{Autoridad única sobre \texttt{/cmd\_vel}}
    Solo un nodo ROS~2 publica \texttt{geometry\_msgs/msg/TwistStamped} en
    \texttt{/cmd\_vel} (verificado previamente en el arranque). Cualquier fallo
    revoca el permiso de movimiento a $(0,0)$.
  \end{block}
\end{frame}

% -------------------------------------------------------------------------
\begin{frame}{6. Modelos y algoritmos empleados}
  \footnotesize
  \begin{tabularx}{\textwidth}{@{}P{3.1cm}P{4.5cm}Y@{}}
    \hline
    \textbf{Componente} & \textbf{Modelo / algoritmo} & \textbf{Función en el sistema}\\
    \hline
    Detección facial &
    \textbf{YuNet} (OpenCV DNN, ONNX) &
    Localización de rostro y 5 puntos clave para alineación afín\\
    Reconocimiento facial &
    \textbf{SFace} ($128$-D) + similitud coseno &
    Descriptores L2 normalizados y decisión dual $(\tau,\mu)$\\
    Detección de mano &
    \textbf{MediaPipe Hand Landmarker} + reglas geométricas &
    21 puntos clave y clasificación angular/distancias en píxeles\\
    Detección de flechas &
    \textbf{Canny} + ajuste de elipses + \textbf{HSV} + geometría de puntas &
    Identificación de discos completos azules/rojos (\texttt{IZQ}/\texttt{DER})\\
    Navegación y control &
    \textbf{Sectores LiDAR por cobertura} + lazo \textbf{yaw} odométrico &
    Evasión conservadora y giros acotados de $90^\circ \pm 5^\circ$\\
    \hline
  \end{tabularx}
  \vspace{0.25em}
  \small
  \begin{block}{Referencia arquitectónica evaluada}
    Se evaluó además una rama alternativa con \texttt{slam\_toolbox} + Nav2 +
    exploración por fronteras antes de priorizar el enfoque reactivo con señales.
  \end{block}
\end{frame}

% -------------------------------------------------------------------------
\begin{frame}{7. Decisiones de diseño y compromisos de ingeniería}
  \footnotesize
  \begin{tabularx}{\textwidth}{@{}P{3.2cm}P{3.6cm}Y@{}}
    \hline
    \textbf{Decisión} & \textbf{Alternativa} & \textbf{Justificación (\textit{trade-off})}\\
    \hline
    YuNet + SFace en CPU &
    GPU / nube / entrenamiento propio &
    Funciona sin internet ni GPU en el lab; depende de iluminación y calidad de galería.\\
    Rechazo $(\tau,\mu)$ (\texttt{NO\_CONCLUYENTE}) &
    1-NN que siempre predice un nombre &
    Evita penalizaciones por identidad falsa en pista, a costa de rechazar tomas borrosas.\\
    Visión geométrica para flechas &
    YOLO reentrenado &
    Cero etiquetado previo y control por tramo; sensible a reflejos o oclusión parcial.\\
    Reactivo + señales vs. SLAM/Nav2 &
    Dependencia exclusiva de \texttt{slam\_toolbox} + Nav2 &
    Menor fragilidad ante desfases de reloj Create~3/Pi y pasillos estrechos; exige integrar meta y vuelta explícitas.\\
    Cobertura LiDAR $\ge 30\%$ &
    Asumir rayos $+\infty$ como libres &
    Evita choques por pérdida de retorno; se ajustó de $50\%$ a $30\%$ tras medir $37$--$47\%$ en hardware.\\
    \hline
  \end{tabularx}
\end{frame}

% -------------------------------------------------------------------------
\begin{frame}{8. Resultados experimentales: facial y pruebas}
  \small
  \begin{columns}[T,onlytextwidth]
    \column{0.535\textwidth}
    \begin{block}{Evaluación facial en particiones reservadas}
      \scriptsize
      \setlength{\tabcolsep}{3pt}
      \begin{tabularx}{\textwidth}{@{}P{2.4cm}cccY@{}}
        \hline
        \textbf{Corte} & \textbf{Cons.} & \textbf{OK} & \textbf{Err.} & \textbf{Neg. sim.}\\
        \hline
        Inicial (26 fotos, 5 id.) & 5 & 4 & 0 & 5/5 rech. (1 rech.)\\
        Intermedio (6 id.) & 6 & 4 & 0 & 6/6 rech. (2 rech.)\\
        Ampliada (74 fotos, 13 id.) & 13 & \textbf{13} & \textbf{0} & \textbf{13/13 rech.}\\
        \hline
      \end{tabularx}
    \end{block}
    \vspace{0.35em}
    \begin{block}{Pruebas unitarias automatizadas}
      \footnotesize
      \begin{itemize}
        \item \textbf{164 pruebas} documentadas (113 Task~1 + 51 Task~2).
        \item \textbf{17 pruebas de señales}: 10 páginas del PDF dan los 4 giros derechos esperados; reflejos validados.
      \end{itemize}
    \end{block}

    \column{0.405\textwidth}
    \begin{block}{Validación en simulación y robot}
      \footnotesize
      \begin{itemize}
        \item \textbf{Task B:} gestos confirmados en Gazebo y hardware + video a 5~FPS.
        \item \textbf{Task A:} evasión + reconocimiento probados juntos en robot real (retornos finitos $37$--$47\%$, pausas hasta $1.66\ \mathrm{s}$).
      \end{itemize}
    \end{block}
    \vspace{0.35em}
    \begin{alertblock}{Resultados de competición}
      \scriptsize [Completar el día del concurso: llegadas A/B, identidad, palabra QR, tiempos, colisiones y puntaje].
    \end{alertblock}
  \end{columns}
\end{frame}

% -------------------------------------------------------------------------
\begin{frame}{9. Lecciones aprendidas y mejoras futuras}
  \small
  \begin{columns}[T,onlytextwidth]
    \column{0.48\textwidth}
    \begin{block}{Lecciones clave del desarrollo}
      \begin{itemize}
        \item \textbf{Dato ausente $\neq$ vía libre:} retornos $+\infty$ del RPLIDAR real obligaron a exigir cobertura $\ge 30\%$ y parada ante TF ausente.
        \item \textbf{Red sin internet en lab:} sincronizar reloj de Pi, fijar \texttt{ROS\_STATIC\_PEERS} y verificar \texttt{TwistStamped} evita fallos silenciosos.
        \item \textbf{Aislar canales:} video a $5\ \mathrm{FPS}$ separado del canal de órdenes protege el frenado.
        \item \textbf{Pruebas por capas:} validar sensores y visión por separado acelera el diagnóstico en hardware.
      \end{itemize}
    \end{block}

    \column{0.46\textwidth}
    \begin{block}{Mejoras futuras}
      \begin{itemize}
        \item Fusionar nube de profundidad OAK-D para detectar \textbf{obstáculos bajos} fuera del plano LiDAR.
        \item Certificar la \textbf{huella circular completa} del chasis frente a choques diagonales.
        \item Completar secuencia de meta A/B, barrido en B, lectura QR exclusiva y retorno automático.
        \item Evaluar impresiones físicas con cámara del robot y migrar parada crítica a bordo de la Pi.
      \end{itemize}
    \end{block}
  \end{columns}
\end{frame}

\end{document}
